\documentclass[a4paper,12pt]{article}

\usepackage[T2A]{fontenc}
\usepackage[utf8]{inputenc}
\usepackage[russian]{babel}

\usepackage{amsmath, amssymb, amsthm, mathtools}
\usepackage{helvet}
\renewcommand{\familydefault}{\sfdefault}
\usepackage{icomma}
\usepackage{geometry}
\usepackage{microtype}
\usepackage{tikz}
\usepackage{tkz-euclide}
\usepackage{pgfplots}
\pgfplotsset{compat=1.18}
\usepackage{tabularx, booktabs, longtable}
\usepackage{enumitem}
\usepackage{hyperref}
\usepackage{parskip}
\usepackage{fancyhdr}
\usepackage{setspace}
\usepackage{xcolor}

\geometry{left=18mm,right=18mm,top=18mm,bottom=20mm}
\onehalfspacing
\sloppy

\definecolor{accent}{HTML}{008080}
\definecolor{darkaccent}{HTML}{004D4D}
\definecolor{textgray}{HTML}{333333}
\definecolor{softgray}{HTML}{F4F6F6}

\hypersetup{
    colorlinks=true,
    linkcolor=darkaccent,
    urlcolor=darkaccent
}

\pagestyle{fancy}
\fancyhf{}
\lhead{\textcolor{textgray}{Чек-лист: задачи на движение}}
\rhead{\textcolor{textgray}{ОГЭ}}
\cfoot{\thepage}
\renewcommand{\headrulewidth}{0.4pt}
\renewcommand{\headrule}{\hbox to\headwidth{\color{accent}\leaders\hrule height \headrulewidth\hfill}}

\setlist[itemize]{leftmargin=*, itemsep=2pt, topsep=3pt}
\setlist[enumerate]{leftmargin=*, itemsep=4pt, topsep=4pt}

\newcommand{\important}[1]{\textcolor{darkaccent}{\textbf{#1}}}
\newcommand{\kmh}{\text{км/ч}}
\newcommand{\km}{\text{км}}
\newcommand{\ch}{\text{ч}}
\newcommand{\minut}{\text{мин}}
\newcommand{\m}{\text{м}}

\newcommand{\simpleline}[5]{
\begin{center}
\begin{tikzpicture}[scale=0.95]
    \draw[line width=0.8pt] (0,0) -- (9,0);
    \fill[accent] (0,0) circle (2pt);
    \fill[accent] (9,0) circle (2pt);
    \node[below] at (0,0) {#1};
    \node[below] at (9,0) {#2};
    \draw[->, line width=1pt, color=darkaccent] (0.5,0.45) -- (4.1,0.45);
    \draw[->, line width=1pt, color=darkaccent] (0.5,0.85) -- (6.0,0.85);
    \node[above] at (2.3,0.45) {#3};
    \node[above] at (3.3,0.85) {#4};
    \node[above] at (4.5,0.05) {#5};
\end{tikzpicture}
\end{center}
}

\newcommand{\circletrack}{
\begin{center}
\begin{tikzpicture}[scale=0.92]
    \draw[line width=1pt, color=accent] (0,0) circle (1.8);
    \fill[darkaccent] (1.8,0) circle (2pt);
    \node[right] at (1.8,0) {старт};
    \draw[->, line width=1pt, color=darkaccent] (1.8,0) arc (0:70:1.8);
    \draw[->, line width=1pt, color=darkaccent] (1.8,0) arc (0:35:1.8);
    \node[above] at (0.5,1.75) {быстрый};
    \node[right] at (1.7,0.8) {медленный};
\end{tikzpicture}
\end{center}
}

\newcommand{\oppositecircletrack}{
\begin{center}
\begin{tikzpicture}[scale=0.92]
    \draw[line width=1pt, color=accent] (0,0) circle (1.8);
    \fill[darkaccent] (1.8,0) circle (2pt);
    \fill[darkaccent] (-1.8,0) circle (2pt);
    \node[right] at (1.8,0) {быстрый};
    \node[left] at (-1.8,0) {медленный};
    \draw[->, line width=1pt, color=darkaccent] (1.8,0) arc (0:130:1.8);
    \draw[->, line width=1pt, color=darkaccent] (-1.8,0) arc (180:230:1.8);
    \node[above] at (0,1.95) {разность пути --- полкруга};
\end{tikzpicture}
\end{center}
}

\begin{document}

\begin{titlepage}
\thispagestyle{empty}
\begin{center}
\vspace*{25mm}

{\Large \textcolor{accent}{\textbf{Чек-лист}}}

\vspace{8mm}

{\Huge \textbf{Текстовые задачи на движение}}

\vspace{4mm}

{\Large разность времён, системы, круговая трасса, движение по воде}

\vspace{22mm}

{\large Лёвин Артём Александрович\\[1mm]
эксклюзивно для Софии Кальней}

\vspace{12mm}

{\large \today}

\vfill

\begin{tikzpicture}
    \draw[color=accent, line width=1.2pt]
    (0,0) .. controls (2,0.8) and (4,-0.8) .. (6,0)
          .. controls (8,0.8) and (10,-0.8) .. (12,0);
    \draw[color=darkaccent, line width=0.5pt]
    (1,-0.35) .. controls (3,0.25) and (5,-0.95) .. (7,-0.35)
          .. controls (9,0.25) and (10.5,-0.75) .. (11.5,-0.35);
\end{tikzpicture}

\vspace{8mm}
\end{center}
\end{titlepage}

\tableofcontents
\newpage

\section{Главная схема решения}

В задачах на движение почти всегда используется одна формула:
\[
S=vt.
\]

Из неё также следуют:
\[
v=\frac{S}{t}, \qquad t=\frac{S}{v}.
\]

\important{Главный рабочий инструмент} --- таблица:
\[
v \quad t \quad S.
\]

Обычно сначала заполняют известные величины, затем обозначают неизвестную за $x$, а третий столбец заполняют по формуле.

\begin{center}
\renewcommand{\arraystretch}{1.35}
\begin{tabularx}{0.86\linewidth}{X c c c}
\toprule
\textbf{Участник движения} & \textbf{Скорость} & \textbf{Время} & \textbf{Расстояние} \\
\midrule
первый & $v_1$ & $t_1$ & $v_1t_1$ \\
второй & $v_2$ & $t_2$ & $v_2t_2$ \\
\bottomrule
\end{tabularx}
\end{center}

\important{Уравнение составляют} по тому столбцу, где появились дроби или выражения, и по той информации из условия, которая ещё не была использована.

\section{Разность времён: кто раньше выехал, тот дольше ехал}

\subsection{Основная идея}

Если два объекта проходят один и тот же путь, но выезжают не одновременно и прибывают одновременно, то тот, кто выехал раньше, находился в пути дольше.

\[
t_{\text{раньше}} - t_{\text{позже}} = \Delta t.
\]

\important{Правило:} если известна разность времён, от большего времени вычитают меньшее.

\[
\text{большее время} - \text{меньшее время} = \text{разность времён}.
\]

\subsection{Пример: грузовик и автобус}

Грузовик и автобус проехали от пункта $A$ до пункта $B$ расстояние $20\ \km$. Грузовик выехал раньше автобуса на $8$ минут, а прибыли они одновременно. Скорость автобуса на $5\ \kmh$ больше скорости грузовика. Найдите скорость автобуса.

\simpleline{$A$}{$B$}{грузовик}{$автобус}{$20\ \km$}

Пусть
\[
x\ \kmh
\]
--- скорость автобуса. Тогда скорость грузовика:
\[
x-5.
\]

\begin{center}
\renewcommand{\arraystretch}{1.35}
\begin{tabularx}{0.82\linewidth}{X c c c}
\toprule
 & $v$ & $t$ & $S$ \\
\midrule
грузовик & $x-5$ & $\dfrac{20}{x-5}$ & $20$ \\
автобус & $x$ & $\dfrac{20}{x}$ & $20$ \\
\bottomrule
\end{tabularx}
\end{center}

Грузовик выехал раньше, значит, он ехал дольше:
\[
\frac{20}{x-5}-\frac{20}{x}=\frac{8}{60}.
\]

\important{ОДЗ:}
\[
x>5.
\]

Решим уравнение:
\[
\frac{20}{x-5}-\frac{20}{x}=\frac{2}{15}.
\]
\[
\frac{20x-20(x-5)}{x(x-5)}=\frac{2}{15}.
\]
\[
\frac{100}{x(x-5)}=\frac{2}{15}.
\]
\[
2x(x-5)=1500.
\]
\[
x(x-5)=750.
\]
\[
x^2-5x-750=0.
\]
\[
D=25+3000=3025=55^2.
\]
\[
x=\frac{5+55}{2}=30.
\]

\[
\boxed{30\ \kmh}
\]

\section{Разность времён при движении в одном направлении}

\subsection{Модель}

Если медленный объект проходит путь дольше, чем быстрый, то:
\[
t_{\text{медленного}}-t_{\text{быстрого}}=\Delta t.
\]

Если быстрый объект движется со скоростью на $a\ \kmh$ больше, чем медленный, то удобно обозначить:
\[
v_{\text{медл}}=x, \qquad v_{\text{быстр}}=x+a.
\]

\subsection{Пример: велосипедист и мотоциклист}

Велосипедист проехал $120\ \km$. Мотоциклист проехал тот же путь со скоростью на $40\ \kmh$ больше и затратил на дорогу на $4$ часа меньше. Найдите скорость велосипедиста.

\begin{center}
\renewcommand{\arraystretch}{1.35}
\begin{tabularx}{0.82\linewidth}{X c c c}
\toprule
 & $v$ & $t$ & $S$ \\
\midrule
велосипедист & $x$ & $\dfrac{120}{x}$ & $120$ \\
мотоциклист & $x+40$ & $\dfrac{120}{x+40}$ & $120$ \\
\bottomrule
\end{tabularx}
\end{center}

Велосипедист движется медленнее, значит, его время больше:
\[
\frac{120}{x}-\frac{120}{x+40}=4.
\]

\important{ОДЗ:}
\[
x>0.
\]

Решаем:
\[
\frac{120(x+40)-120x}{x(x+40)}=4.
\]
\[
\frac{4800}{x(x+40)}=4.
\]
\[
x(x+40)=1200.
\]
\[
x^2+40x-1200=0.
\]
\[
D=1600+4800=6400=80^2.
\]
\[
x=\frac{-40+80}{2}=20.
\]

\[
\boxed{20\ \kmh}
\]

\section{Две ситуации и система уравнений}

\subsection{Когда появляется система}

Система уравнений обычно появляется, если в задаче есть две разные ситуации движения.

Например:

\begin{enumerate}
    \item сначала два объекта едут к одной точке;
    \item затем из этой точки движутся в разные стороны.
\end{enumerate}

В таком случае удобно составить две таблицы и получить два уравнения.

\subsection{Пример с пунктами $A$, $C$, $B$}

Автомобиль выехал из пункта $A$ в пункт $C$. Через $15$ минут после него из пункта $A$ в пункт $C$ выехал мотоциклист со скоростью $75\ \kmh$. В пункт $C$ они прибыли одновременно. Затем из пункта $C$ автомобиль поехал в пункт $B$, а мотоциклист --- обратно в сторону пункта $A$. К моменту прибытия автомобиля в пункт $B$ мотоциклист проехал $\dfrac{2}{3}$ пути от $C$ до $A$. Расстояние $AB$ равно $115\ \km$.

Пусть:
\[
AC=x, \qquad v_{\text{авт}}=v.
\]

\begin{center}
\begin{tikzpicture}[scale=0.95]
    \draw[line width=0.8pt] (0,0) -- (9,0);
    \fill[accent] (0,0) circle (2pt);
    \fill[accent] (5.2,0) circle (2pt);
    \fill[accent] (9,0) circle (2pt);
    \node[below] at (0,0) {$A$};
    \node[below] at (5.2,0) {$C$};
    \node[below] at (9,0) {$B$};
    \draw[<->, color=darkaccent] (0,0.55) -- (5.2,0.55);
    \node[above] at (2.6,0.55) {$x$};
    \draw[<->, color=darkaccent] (5.2,0.55) -- (9,0.55);
    \node[above] at (7.1,0.55) {$115-x$};
\end{tikzpicture}
\end{center}

\subsection*{Первая ситуация: движение из $A$ в $C$}

\begin{center}
\renewcommand{\arraystretch}{1.35}
\begin{tabularx}{0.82\linewidth}{X c c c}
\toprule
 & $v$ & $t$ & $S$ \\
\midrule
автомобиль & $v$ & $\dfrac{x}{v}$ & $x$ \\
мотоциклист & $75$ & $\dfrac{x}{75}$ & $x$ \\
\bottomrule
\end{tabularx}
\end{center}

Автомобиль выехал на $15$ минут раньше, значит, он ехал дольше:
\[
\frac{x}{v}-\frac{x}{75}=\frac{15}{60}.
\]

\subsection*{Вторая ситуация: движение из $C$}

Автомобиль проходит:
\[
CB=115-x.
\]

Мотоциклист проходит:
\[
\frac{2}{3}x.
\]

\begin{center}
\renewcommand{\arraystretch}{1.35}
\begin{tabularx}{0.82\linewidth}{X c c c}
\toprule
 & $v$ & $t$ & $S$ \\
\midrule
автомобиль & $v$ & $\dfrac{115-x}{v}$ & $115-x$ \\
мотоциклист & $75$ & $\dfrac{\frac{2}{3}x}{75}$ & $\dfrac{2}{3}x$ \\
\bottomrule
\end{tabularx}
\end{center}

Во второй ситуации они движутся одновременно, поэтому времена равны:
\[
\frac{115-x}{v}=\frac{\frac{2}{3}x}{75}.
\]

Итоговая математическая модель:
\[
\begin{cases}
\dfrac{x}{v}-\dfrac{x}{75}=\dfrac{15}{60},\\[3mm]
\dfrac{115-x}{v}=\dfrac{\frac{2}{3}x}{75}.
\end{cases}
\]

\important{Метод решения:} чаще всего удобно выразить одну переменную из второго уравнения и подставить в первое.

\section{Движение по круговой трассе}

\subsection{Старт из одной точки в одном направлении}

Если два объекта стартуют из одной точки в одном направлении и быстрый впервые обгоняет медленного на круг, то быстрый проходит на один круг больше.

\circletrack

Если длина круга равна $L$, то:
\[
S_{\text{быстр}}-S_{\text{медл}}=L.
\]

Если время движения равно $t$, то:
\[
v_{\text{быстр}}t-v_{\text{медл}}t=L.
\]

То есть:
\[
(v_{\text{быстр}}-v_{\text{медл}})t=L.
\]

\subsection{Пример}

Два автомобиля одновременно стартуют из одной точки круговой трассы длиной $15\ \km$ и движутся в одном направлении. Скорости автомобилей равны $80\ \kmh$ и $60\ \kmh$. Через сколько минут быстрый автомобиль впервые обгонит медленный на круг?

\begin{center}
\renewcommand{\arraystretch}{1.35}
\begin{tabularx}{0.72\linewidth}{X c c c}
\toprule
 & $v$ & $t$ & $S$ \\
\midrule
быстрый & $80$ & $x$ & $80x$ \\
медленный & $60$ & $x$ & $60x$ \\
\bottomrule
\end{tabularx}
\end{center}

Быстрый проходит на один круг больше:
\[
80x-60x=15.
\]
\[
20x=15.
\]
\[
x=\frac{15}{20}=\frac{3}{4}\ \ch.
\]

Переводим в минуты:
\[
\frac{3}{4}\cdot 60=45.
\]

\[
\boxed{45\ \minut}
\]

\subsection{Найти скорость медленного}

Два автомобиля стартовали одновременно из одной точки круговой трассы длиной $10\ \km$ и двигались в одном направлении. Скорость быстрого автомобиля равна $90\ \kmh$. Через $40$ минут он впервые обогнал медленного на круг. Найдите скорость медленного автомобиля.

Переводим время:
\[
40\ \minut=\frac{40}{60}=\frac{2}{3}\ \ch.
\]

Пусть $x$ --- скорость медленного автомобиля.

\begin{center}
\renewcommand{\arraystretch}{1.35}
\begin{tabularx}{0.72\linewidth}{X c c c}
\toprule
 & $v$ & $t$ & $S$ \\
\midrule
быстрый & $90$ & $\dfrac{2}{3}$ & $90\cdot\dfrac{2}{3}$ \\
медленный & $x$ & $\dfrac{2}{3}$ & $x\cdot\dfrac{2}{3}$ \\
\bottomrule
\end{tabularx}
\end{center}

\[
90\cdot\frac{2}{3}-x\cdot\frac{2}{3}=10.
\]
\[
60-\frac{2x}{3}=10.
\]
\[
\frac{2x}{3}=50.
\]
\[
x=75.
\]

\[
\boxed{75\ \kmh}
\]

\section{Круговая трасса: старт из противоположных точек}

Если два объекта стартуют из диаметрально противоположных точек и движутся в одном направлении, то между ними сначала половина круга.

\oppositecircletrack

Поэтому при первом догоне быстрый проходит больше медленного не на целый круг, а на половину круга:
\[
S_{\text{быстр}}-S_{\text{медл}}=\frac{L}{2}.
\]

\subsection{Пример}

Два мотоциклиста стартовали одновременно из диаметрально противоположных точек круговой трассы длиной $20\ \km$ и поехали в одном направлении. Скорость быстрого мотоциклиста на $12\ \kmh$ больше скорости медленного. Через сколько минут быстрый мотоциклист догонит медленного?

Пусть:
\[
v_{\text{медл}}=y, \qquad v_{\text{быстр}}=y+12.
\]

\begin{center}
\renewcommand{\arraystretch}{1.35}
\begin{tabularx}{0.78\linewidth}{X c c c}
\toprule
 & $v$ & $t$ & $S$ \\
\midrule
быстрый & $y+12$ & $x$ & $x(y+12)$ \\
медленный & $y$ & $x$ & $xy$ \\
\bottomrule
\end{tabularx}
\end{center}

Разность путей равна половине круга:
\[
x(y+12)-xy=\frac{20}{2}.
\]
\[
xy+12x-xy=10.
\]
\[
12x=10.
\]
\[
x=\frac{10}{12}=\frac{5}{6}\ \ch.
\]

Переводим в минуты:
\[
\frac{5}{6}\cdot 60=50.
\]

\[
\boxed{50\ \minut}
\]

\section{Движение по воде}

\subsection{Собственная скорость, течение, плот}

Пусть:
\[
x \text{ --- собственная скорость лодки или теплохода},
\]
\[
u \text{ --- скорость течения}.
\]

Тогда:

\begin{center}
\renewcommand{\arraystretch}{1.35}
\begin{tabularx}{0.82\linewidth}{X c}
\toprule
\textbf{Ситуация} & \textbf{Скорость} \\
\midrule
по течению & $x+u$ \\
против течения & $x-u$ \\
в неподвижной воде & $x$ \\
плот по течению & $u$ \\
\bottomrule
\end{tabularx}
\end{center}

\important{ОДЗ для движения против течения:}
\[
x>u.
\]

\subsection{Пример: путь туда и обратно}

Баржа прошла против течения $24\ \km$ и вернулась обратно. На обратный путь она затратила на $3$ часа меньше, чем на путь против течения. Скорость течения равна $2\ \kmh$. Найдите собственную скорость баржи.

Пусть $x$ --- собственная скорость баржи.

\begin{center}
\renewcommand{\arraystretch}{1.35}
\begin{tabularx}{0.82\linewidth}{X c c c}
\toprule
 & $v$ & $t$ & $S$ \\
\midrule
против течения & $x-2$ & $\dfrac{24}{x-2}$ & $24$ \\
по течению & $x+2$ & $\dfrac{24}{x+2}$ & $24$ \\
\bottomrule
\end{tabularx}
\end{center}

Против течения баржа движется медленнее, значит, времени тратит больше:
\[
\frac{24}{x-2}-\frac{24}{x+2}=3.
\]

\important{ОДЗ:}
\[
x>2.
\]

Решаем:
\[
\frac{24(x+2)-24(x-2)}{(x-2)(x+2)}=3.
\]
\[
\frac{96}{x^2-4}=3.
\]
\[
3x^2-12=96.
\]
\[
3x^2=108.
\]
\[
x^2=36.
\]
\[
x=6.
\]

\[
\boxed{6\ \kmh}
\]

\subsection{Пример со стоянкой}

Теплоход прошёл некоторое расстояние по течению, сделал стоянку на $4$ часа и вернулся обратно. Его собственная скорость равна $20\ \kmh$, скорость течения равна $4\ \kmh$. В исходный пункт теплоход вернулся через $14$ часов после отплытия. Найдите общее расстояние, пройденное теплоходом.

Скорость по течению:
\[
20+4=24\ \kmh.
\]

Скорость против течения:
\[
20-4=16\ \kmh.
\]

Пусть $x$ --- расстояние в одну сторону.

\begin{center}
\renewcommand{\arraystretch}{1.35}
\begin{tabularx}{0.82\linewidth}{X c c c}
\toprule
 & $v$ & $t$ & $S$ \\
\midrule
по течению & $24$ & $\dfrac{x}{24}$ & $x$ \\
стоянка & $0$ & $4$ & $0$ \\
против течения & $16$ & $\dfrac{x}{16}$ & $x$ \\
\bottomrule
\end{tabularx}
\end{center}

Общее время равно $14$ часам:
\[
\frac{x}{24}+4+\frac{x}{16}=14.
\]
\[
\frac{x}{24}+\frac{x}{16}=10.
\]

Приводим к общему знаменателю:
\[
\frac{2x}{48}+\frac{3x}{48}=10.
\]
\[
\frac{5x}{48}=10.
\]
\[
x=96.
\]

Общее расстояние:
\[
2x=192.
\]

\[
\boxed{192\ \km}
\]

\subsection{Пример с плотом}

Лодка прошла $45\ \km$ по течению и $45\ \km$ против течения. Плот прошёл $32\ \km$ по течению. Плот отправился на $1$ час раньше лодки, а прибыли они одновременно. Скорость течения равна $4\ \kmh$. Найдите собственную скорость лодки.

Пусть $x$ --- собственная скорость лодки.

\begin{center}
\renewcommand{\arraystretch}{1.35}
\begin{tabularx}{0.86\linewidth}{X c c c}
\toprule
 & $v$ & $t$ & $S$ \\
\midrule
лодка по течению & $x+4$ & $\dfrac{45}{x+4}$ & $45$ \\
лодка против течения & $x-4$ & $\dfrac{45}{x-4}$ & $45$ \\
плот по течению & $4$ & $\dfrac{32}{4}=8$ & $32$ \\
\bottomrule
\end{tabularx}
\end{center}

Плот отправился раньше, значит, он двигался дольше:
\[
8-\left(\frac{45}{x+4}+\frac{45}{x-4}\right)=1.
\]

То есть:
\[
\frac{45}{x+4}+\frac{45}{x-4}=7.
\]

\important{ОДЗ:}
\[
x>4.
\]

Решаем:
\[
\frac{45(x-4)+45(x+4)}{x^2-16}=7.
\]
\[
\frac{90x}{x^2-16}=7.
\]
\[
90x=7x^2-112.
\]
\[
7x^2-90x-112=0.
\]
\[
D=8100+3136=11236=106^2.
\]
\[
x=\frac{90+106}{14}=14.
\]

\[
\boxed{14\ \kmh}
\]

\section{Алгоритм решения}

\begin{enumerate}
    \item Определите тип движения: по прямой, по кругу, по воде, с задержкой старта, в две ситуации.
    \item Составьте таблицу $v$, $t$, $S$.
    \item Введите переменную именно для той величины, которую удобно найти первой.
    \item Заполните третий столбец по формуле $S=vt$ или $t=\dfrac{S}{v}$.
    \item Используйте информацию из условия, которая ещё не вошла в таблицу.
    \item Если дана разность времён, от большего времени вычитайте меньшее.
    \item Если движение по кругу, записывайте разность путей: один круг или половина круга.
    \item Если движение по воде, отдельно учитывайте течение: $x+u$, $x-u$, $u$.
    \item Если есть две ситуации, составляйте две таблицы и систему.
    \item После решения проверьте, что ответ дан в нужных единицах.
\end{enumerate}

\section{Типичные ошибки}

\begin{enumerate}
    \item \important{Не перевести минуты в часы.} Если скорость в $\kmh$, то время должно быть в часах.
    \item \important{Вычесть времена в неправильном порядке.} Нужно вычитать меньшее время из большего.
    \item \important{Перепутать скорость по течению и против течения.} По течению: $x+u$, против течения: $x-u$.
    \item \important{Забыть, что плот движется со скоростью течения.} У плота нет собственной скорости.
    \item \important{На круговой трассе взять целый круг вместо половины.} Если старт из противоположных точек, разность путей равна половине круга.
    \item \important{Смешать разные неизвестные.} Если $x$ уже обозначает расстояние, не нужно тем же $x$ обозначать скорость.
    \item \important{Ответить на промежуточную величину.} Если найдено расстояние в одну сторону, а спрашивают общий путь, нужно умножить на $2$.
    \item \important{Потерять ОДЗ.} В задачах с дробями знаменатели не должны обращаться в ноль; в движении против течения собственная скорость должна быть больше скорости течения.
\end{enumerate}

\newpage

\section{Тренировочный блок}

Решите задачи. Ответы приведены на отдельной странице после тренировочного блока.

\subsection*{Средний уровень}

\begin{enumerate}
    \item Грузовик и автобус проехали от пункта $A$ до пункта $B$ расстояние $30\ \km$. Грузовик выехал на $10$ минут раньше автобуса, а прибыли они одновременно. Скорость автобуса на $6\ \kmh$ больше скорости грузовика. Найдите скорость автобуса.

    \item Велосипедист проехал $120\ \km$. Мотоциклист проехал тот же путь со скоростью на $40\ \kmh$ больше и затратил на дорогу на $4$ часа меньше. Найдите скорость велосипедиста.

    \item Два автомобиля одновременно стартуют из одной точки круговой трассы длиной $18\ \km$ и движутся в одном направлении. Их скорости равны $75\ \kmh$ и $60\ \kmh$. Через сколько минут быстрый автомобиль впервые обгонит медленный на круг?
\end{enumerate}

\subsection*{Выше среднего уровня}

\begin{enumerate}[resume]
    \item Два автомобиля стартовали одновременно из одной точки круговой трассы длиной $12\ \km$ и двигались в одном направлении. Скорость быстрого автомобиля равна $90\ \kmh$. Через $36$ минут он впервые обогнал медленный автомобиль на круг. Найдите скорость медленного автомобиля.

    \item Два мотоциклиста стартовали одновременно из диаметрально противоположных точек круговой трассы длиной $36\ \km$ и поехали в одном направлении. Скорость быстрого мотоциклиста на $18\ \kmh$ больше скорости медленного. Через сколько часов быстрый мотоциклист догонит медленного?

    \item Катер прошёл против течения $24\ \km$ и вернулся обратно. На обратный путь он затратил на $3$ часа меньше, чем на путь против течения. Скорость течения равна $2\ \kmh$. Найдите собственную скорость катера.

    \item Теплоход прошёл некоторое расстояние по течению, сделал стоянку на $2$ часа и вернулся обратно. Его собственная скорость равна $18\ \kmh$, скорость течения равна $3\ \kmh$. В исходный пункт теплоход вернулся через $8$ часов после отплытия. Найдите общее расстояние, пройденное теплоходом.

    \item Лодка прошла $48\ \km$ по течению и $48\ \km$ против течения, затратив на весь путь $9$ часов. Скорость течения равна $4\ \kmh$. Найдите собственную скорость лодки.

    \item Автомобиль выехал из пункта $A$ в пункт $C$. Через $18$ минут после него из пункта $A$ в пункт $C$ выехал мотоциклист со скоростью $75\ \kmh$. В пункт $C$ они прибыли одновременно. Затем из пункта $C$ автомобиль поехал в пункт $B$, а мотоциклист --- обратно в сторону пункта $A$. К моменту прибытия автомобиля в пункт $B$ мотоциклист проехал $\dfrac{2}{3}$ пути от $C$ до $A$. Расстояние $AB$ равно $65\ \km$. Найдите расстояние $AC$.
\end{enumerate}

\subsection*{Сложный уровень}

\begin{enumerate}[resume]
    \item Лодка прошла $45\ \km$ по течению и $45\ \km$ против течения. Плот прошёл $32\ \km$ по течению. Плот отправился на $1$ час раньше лодки, а прибыли они одновременно. Скорость течения равна $4\ \kmh$. Найдите собственную скорость лодки.
\end{enumerate}

\newpage

\section{Ответы к тренировочному блоку}

\begin{center}
\renewcommand{\arraystretch}{1.35}
\begin{tabularx}{0.9\linewidth}{c X}
\toprule
\textbf{№} & \textbf{Ответ} \\
\midrule
1 & $36\ \kmh$ \\
2 & $20\ \kmh$ \\
3 & $72\ \minut$ \\
4 & $70\ \kmh$ \\
5 & $1\ \ch$ \\
6 & $6\ \kmh$ \\
7 & $105\ \km$ \\
8 & $12\ \kmh$ \\
9 & $45\ \km$ \\
10 & $14\ \kmh$ \\
\bottomrule
\end{tabularx}
\end{center}

\end{document}